\section{Method}\label{sec:method}
The method constructs bounded financial and economic decision games with
explicit actor goals and incentives. It follows the order of execution: construct a world, restrict each
participant's observation, obtain a permitted action, settle the complete
round, and preserve the result. We use \emph{participant} for a modeled
decision maker and \emph{persona} for its descriptive role, motivation, and
goals. These terms describe economic roles and their decision incentives; they do
not imply that the participant represents an identified real person. Table~\ref{tab:notation} fixes the notation used below.

\begin{table}[t]
\caption{Core notation. Each symbol has one role throughout the method.}
\label{tab:notation}
\centering\footnotesize
\begin{tabularx}{\columnwidth}{@{}p{0.23\columnwidth}Y@{}}
\toprule Symbol & Meaning and reason for retaining it \\
\midrule
$i,t,k$ & Participant, round, and metric indices; identify whose decision and which state are being discussed. \\
$b_{i,t}$ & Participant $i$'s cash before decisions in round $t$; limits feasible spending. \\
$x_{k,t}$ & Public metric $k$ before decisions in round $t$; records the shared situation. \\
$A_{i,t},a_{i,t}$ & Affordable menu and selected action; separate permission from choice. \\
$c_i(a)$ & Cash cost of action $a$; connects a choice to resource use. \\
$q_{i,t}(a)$ & Model-returned score for an affordable action; used for explicit selection. \\
$U_i,w_{i,k}$ & Participant utility and its metric weights; state the assumed trade-offs. The cash weight is $w_{i,\mathrm{cash}}$. \\
\bottomrule
\end{tabularx}
\end{table}

\subsection{Constructing the Game Before the Persona}\label{sec:construction}
A persona becomes strategically meaningful through its preferences, feasible
choices, resources, and information. A biography alone does not define a game.
In a game, a participant's preferred outcome can depend on what others choose;
players, actions, and payoffs must therefore be specified together
\cite{polak2,nash1950}. Boyun treats world construction as a research task
performed by a user or separate authoring model, followed by explicit review.
Jev selects actions during execution; it does not write the world or invent
its payoff rules.

Construction begins with the research question, focal participant, horizon,
and outcome of interest. The author then includes roles because they control
a consequential resource, know different information, or can change another
participant's outcome. For example, a reviewer matters if verification changes
credibility or another actor's incentive to participate. Adding a reviewer
whose permitted actions affect nothing that anyone values would create
narrative detail without that strategic mechanism. The construction review
reports objective-sensitive dependencies, identifying action effects that
change another participant's declared utility at the initial bounded state.
These are local diagnostics, not a proof of strategic independence in all
future states.

Each role is turned into a participant with a motivation, terminal goal,
resource allocation, information boundary, utility weights, and action menu.
Motivation explains why the participant values particular outcomes; the goal
states the intended preference; weights and action effects make the trade-off
executable. A participant maximizing its own utility need not maximize money:
a reviewer can value credibility, while a treasury steward values retained
resources and low execution risk. The author must explain the units and a
numerical trade-off for each objective. Schema validation cannot establish
that a stated goal and its supplied weights are semantically aligned; the
review must examine both.

Different role names need not denote different game-theoretic types. Here,
variants differ operationally through preferences, resources, available
actions, private information, or observation links. The specification does
not supply a probability distribution over private types, so these variants
do not by themselves constitute a solved Bayesian game. Similarly, a linear
penalty on an execution-risk metric represents sensitivity to that stipulated
metric; it does not implement formal risk aversion over uncertain wealth.
Describing an actor as cautious or speculative cannot supply the missing
probability and utility model.

The author next fixes timing and conditional consequences. Current actions
are simultaneous; preset messages arrive in the next round; specified
joint-action rules express complementarity or conflict. An action is a
choice at one information state, whereas a strategy is a contingent plan
across relevant information states \cite{polak15,polak18}. Boyun requests
successive actions from observations; it does not solve an extensive-form
game or establish subgame-perfect behavior. Threats, promises, and commitments
have no enforcement beyond the implemented costs and transition rules.
Without that distinction, a model could treat persuasive language as a binding
contract that the engine never enforces.

Every consequential coefficient, motivation, allocation, and link records a
rationale and provenance status. A news source can motivate a role or event
without identifying its numeric effect. The construction review exposes
initial cash allocations, utility units, affordable actions, isolated action
consequences, and conditional interactions. It warns about zero objectives,
a single feasible option, tied isolated utilities, absent initial objective
coupling, or unavailable peer messages. A warning requests interpretation:
a control participant can reasonably have no message links. Passing these
checks establishes a sufficiently explicit model to inspect, not a
psychologically valid population.

\subsection{Cast Size and a Worked Construction}\label{sec:cast}
Cast size follows the mechanism being investigated. Two participants can
isolate a coordination problem; a small heterogeneous cast can expose
conflicting incentives. Ten, twenty, fifty, or one hundred participants are
useful only when additional variants test a stated information or preference
hypothesis. The supported range is 2--100 participants, 1--8 rounds, 1--8
metrics, and 2--8 actions per participant, subject to aggregate size limits.
These bounds do not establish that one hundred synthetic actors represent a
market or provide one hundred independent observations.

Table~\ref{tab:cast} specifies an executable eight-participant example: a
hypothetical NEAR DeFi pilot with two rounds and a total cast budget of
\$10,000. The shared metrics are readiness, credibility, participation, and
execution risk, each bounded from zero to one hundred hypothetical points.
The initial values are 20, 40, 20, and 40 respectively. All allocations,
preferences, and effects are explicit assumptions, not estimated NEAR market
behavior. Spending purchases modeled effort; no token, liquidity-provider
position, asset income, or creator profit is calculated.

\begin{table}[t]
\caption{Assumed pilot cast. Cash allocations sum to \$10,000. Every participant
also has a silent, zero-cost wait option. Goals are implemented by reviewed
utility weights, not inferred from role labels.}
\label{tab:cast}
\centering\footnotesize
\begin{tabularx}{\columnwidth}{@{}p{0.24\columnwidth}rY@{}}
\toprule Role & Cash (\$) & Objective emphasis and active options \\
\midrule
Founder & 3,000 & Readiness and credible participation; fund a pilot or announce broadly. \\
Treasury steward & 2,000 & Runway and low execution risk; publish a resource disclosure. \\
Liquidity sponsor & 1,500 & Readiness with manageable risk; contribute integration effort. \\
Independent reviewer & 200 & Credibility and lower risk; spend effort on review. \\
Risk-sensitive retail & 700 & Credibility and low risk; join the limited trial. \\
Participation-oriented retail & 700 & Participation with a weaker risk penalty; join and share. \\
Large participant & 1,500 & Readiness and participation versus effort and risk; contribute larger trial effort. \\
Ecosystem partner & 400 & Credible participation versus outreach cost; amplify the trial. \\
\bottomrule
\end{tabularx}
\end{table}

Preferences create inspectably different choices even before strategic
responses are considered. For the participation-oriented participant, joining
costs 5,000 cents, adds five participation points, and adds three risk points.
With one utility point per retained cent, 3,000 per participation point, and
minus 500 per risk point, its isolated utility change is
$-5{,}000+3{,}000\times5-500\times3=8{,}500$ points.
The risk-sensitive participant's join action costs the same, adds three
participation points and one risk point; its weights of 100 and minus 3,000
produce $-5{,}000+100\times3-3{,}000=-7{,}700$ points.
These are partial current-stage comparisons with waiting, not terminal
strategy values or estimates of what actual retail users will do.

Joint rules give these isolated comparisons strategic context. In this example,
funding and independent review together add credibility and reduce execution
risk; sponsor effort combined with the risk-sensitive participant joining
adds participation and reduces risk; broad announcement combined with
participation-oriented joining raises both participation and risk. Other
actors can therefore change which outcomes are attractive. Removing these
rules, changing preferences, or changing message links defines separate
experiments rather than cosmetic persona edits.

Increasing persona detail must preserve the intended economics. Splitting one
role into variants requires explicit reallocation of its budget and review of
per-action effects and links; copying its full resources into each variant
would change the game. Equal total cash alone does not preserve aggregate
influence when each extra actor can apply a full metric effect. The current
engine does not automatically infer population weights or perform a
budget-preserving split. Such comparisons require separately reviewed
specifications, a stable evidence cutoff, and a stated reason for each variant.

\subsection{Information and Communication}
An observation contains a participant's own motivation, goal, utility weights,
cash, affordable actions, public metrics, permitted rules, delivered messages,
and visible events. It excludes other participants' private goals and future
event text. This projection is necessary to avoid giving every participant
the author's complete knowledge of the world. The host still possesses that
knowledge; restricting an observation is not encryption.

The directed information graph specifies subscriptions: an edge identifies
whose publications a participant receives. An action's preset publication is
delivered to subscribers in the next round. A message is a claim that may
influence a decision, not proof of its contents. Delaying delivery makes the
timing explicit and prevents an early request in a round from influencing a
later request in that same round. A private event's text can remain restricted
even when its declared effect changes a public metric.

This is a bounded publication mechanism rather than unrestricted conversation.
Participants do not independently call tools or negotiate arbitrary new
contracts. A scenario requiring such behavior would need additional action
types and settlement rules before execution.

\subsection{Affordable Actions and Typed Selection}
Feasibility is determined before inference. For a participant's reviewed menu,
the affordable set is
\begin{equation}
A_{i,t}=\{a\in\operatorname{Menu}_i:c_i(a)\leq b_{i,t}\}.
\label{eq:affordable}
\end{equation}
Equation~\eqref{eq:affordable} makes permission depend on current cash rather
than on the model's interpretation of a budget. Without it, a plausible action
could spend unavailable funds. The engine validates affordability again at
settlement to reject stale or altered decisions.

Jev is the primary model for participant action selection because its native
interface matches a bounded game decision: state plus typed questions produce
answers and scores rather than an arbitrary prose action
\cite{typesafeintro,typesafeapi}. World authoring remains a separate language
model or human task. This separation places contextual action judgment in Jev
and resource accounting, arithmetic, and enforcement in code. It avoids asking
a prose response to serve simultaneously as a decision, an executable action,
and a trustworthy numerical calculation. TypeSafe's architecture and efficiency
claims are provider descriptions; this design choice is not evidence that Jev
reproduces people or outperforms other models in strategic reasoning.

The instruction contract tells the model to act as its own participant and
seek its declared terminal utility, considering the remaining horizon and
visible interactions. Motivation provides context but cannot override costs,
weights, or action permissions. The model must not assume access to opponents'
private preferences, invent event likelihoods, or treat a message as an enforced
commitment. Code supplies each option's exact isolated current-stage cash,
bounded metrics, and utility change, holding other actions fixed. This avoids
outsourcing arithmetic to the model. The projection explicitly excludes peer
choices, joint effects, and later rounds, so its largest delta is not a
prescribed strategy. The instruction expresses a decision objective, not a
guarantee of rational or optimal behavior.

For each non-forced turn, the adapter sends the observation
and a Choice question over $A_{i,t}$. It requires one finite score
$q_{i,t}(a)$ between zero and one for every option, a total score of one, the
maximum-score label, and the response's confidence field. Missing options or
an invalid distribution reject the decision. If only wait is feasible, no
model request is needed.

Selection either takes the highest-scored action or samples one action from
the returned distribution. The chosen mode is fixed before execution. Sampling
permits different paths under the same stated preferences; highest-score
selection gives a direct deterministic interpretation of a saved response.
Local sampling uses a recorded seed and observation identity. Neither method
makes the scores calibrated frequencies of real behavior. In particular,
$q_{i,t}(a)$ is not a probability that the studied intervention succeeds and is not a
mixed-strategy equilibrium. Provider confidence is retained as a response
field, without assuming an established statistical confidence interpretation.

The request includes full action descriptions once in the observation;
answer criteria reference action identifiers and the computed projections.
Removing repeated descriptions reduces duplication for descriptive menus,
although the additional numeric support can enlarge a short request. One Choice
question covers one actor's menu; there is no extra explanatory model call.
Stable menu order, a versioned instruction contract, the exact request identity,
and the actual returned model identity make prompt changes traceable. Whether
this structure improves decision quality or cost at equal quality requires a
controlled comparison; smaller payloads alone do not establish that result.

\subsection{Synchronous Settlement and Replay}\label{sec:settlement}
All participants decide from the same pre-settlement round state. The engine
then applies the complete action profile, rather than mutating the world
after each model request. Otherwise, request order would confer an unintended
information or execution advantage. Algorithm~\ref{alg:round} gives the
procedure without depending on implementation-specific function names.

\begin{algorithm}[t]
\caption{World execution and recorded replay}
\label{alg:round}
\footnotesize
\begin{algorithmic}[1]
\Require Reviewed world, selection mode, recorded seed
\State Validate participants, costs, metrics, events, and links
\For{each round $t$}
  \State Apply this round's scheduled metric changes
  \State Aggregate changes before enforcing metric bounds
  \State Deliver permitted publications from earlier rounds
  \State Freeze the pre-settlement state
  \For{each participant $i$}
    \State Build own observation and $A_{i,t}$
    \If{only wait is feasible}
      \State Record wait without inference
    \Else
      \State Obtain and validate typed action scores
      \State Select $a_{i,t}$ using the reviewed mode
    \EndIf
  \EndFor
  \State Validate every decision against its observation
  \If{any decision is invalid}
    \State Reject the round without partial settlement
    \State \textbf{stop}
  \Else
    \State Stage subtraction of each cost $c_i(a_{i,t})$
    \State Sum action effects and matching joint-rule effects
    \State Enforce metric bounds once on the aggregate
    \State Commit all staged updates together
    \State Queue preset publications for the next round
    \State Save decisions, receipts, and settled state
  \EndIf
\EndFor
\State Replay by replacing score requests with saved decisions
\State Compare recomputed states with recorded states
\end{algorithmic}
\end{algorithm}

An action effect is the metric change declared for one participant's action.
A joint-action rule adds a declared change when its specified combination of
actions occurs, allowing complementarity or conflict. For example, funding
readiness and a partner joining may have a combined effect beyond either
action alone. Without a joint rule, the model assumes their effects are
additive. Whether complementarity exists and how large it is remain world
assumptions, not properties established by the engine.

Metric changes are summed before enforcing their bounds. If readiness is
bounded by ten, starts at nine, and receives changes of plus four and minus
four, aggregate settlement returns nine. Applying the upper bound after the
positive action alone would instead return six after the negative action.
The aggregate rule avoids that dependence on iteration order. Scheduled events
are a separate stage before the next observation; a plotted settlement value
can therefore differ from the following round's opening value.

\begin{proposition}[Resource accounting]
With nonnegative initial cash, affordable actions preserve nonnegative cash.
Remaining cash plus cumulative action costs equals initial cash at every
settled round.
\end{proposition}
\begin{IEEEproof}
Equation~\eqref{eq:affordable} bounds every cost by the participant's available
cash. Settlement subtracts that same cost from cash and adds it to cumulative
spending. Each subtraction therefore preserves the stated total, starting
from the initial balances.
\end{IEEEproof}
This invariant detects unexplained resource creation or loss. Action costs
leave the modeled participants rather than being transferred to an explicit
counterparty, so the proposition concerns a cost-accounting model, not a
complete financial system or a solvency guarantee.

Replay reuses recorded actions and recomputes their consequences. It is
needed because a later model request can return a different action or model
version. Without replay, a numerical discrepancy could be confused with
changed inference. Exact agreement establishes consistency with the recorded
rules and choices, not truth of the rules themselves.

\subsection{What the Decision Record Can Explain}\label{sec:record-scope}
The record makes the experiment inspectable without revealing the model's
internal reasoning. It preserves the reviewed world, each supplied observation,
the typed response, the selected action and selection mode, and the resulting
numeric transitions. An investigator can check affordability, compare the
selected action with a supplied payoff table, inspect alternative scores,
and replay the consequences. These are input/output and rule-execution checks.

Jev returns typed answers and probabilities, not a deliberation transcript
\cite{typesafeapi}. Its confidence field does not explain a preference, and
retained motivations describe the modeled actor rather than an observed reason
for its choice. A code-computed utility comparison can show that an action
would improve a specified payoff; it cannot establish that Jev used that
comparison internally. A later prose explanation from a separate model would
be another generated interpretation, not recovered evidence of Jev's process.

Exact replay recomputes settlement using the saved choices; it does not
reproduce the hidden inference or guarantee the same answer on a new request.
A provider response reports its model identity without the present adapter
verifying a model-execution attestation. A NEAR commitment can bind disclosed
artifact bytes to a registered fingerprint, but cannot establish why an action
was chosen or whether its incentives describe real actors. We therefore claim
explicit assumptions, recorded choices, and replayable numerical execution,
rather than reasoning transparency or comprehensive auditability.

\subsection{Utility and Game Diagnostics}
Utility expresses how a participant values outcomes; it prevents ``better''
from meaning the same thing for every role. The current model uses
\begin{equation}
U_i=w_{i,\mathrm{cash}}b_{i,t}+\sum_k w_{i,k}x_{k,t}.
\label{eq:utility}
\end{equation}
The weights convert cash and metrics into a common assumed preference scale.
For instance, a participant may trade some cash for readiness while another
values retained cash more strongly. Equation~\eqref{eq:utility} makes that
trade-off inspectable. Without explicit weights, a reported improvement has
no defined meaning for the participant. The weights and units must be reviewed;
they are not inferred psychological traits, and utility points are not profit.
The same expression is used for settled and terminal states by substituting
their recorded balances and metrics.

A separate small-game diagnostic enumerates joint action profiles, computes
their utilities, and checks unilateral deviations. It reports best responses,
pure-strategy Nash equilibria, and actions that are strictly inferior to another
pure action. Given supplied beliefs about opponents, it computes expected
utility as a probability-weighted average; regret is the gap between a chosen
action and the best available expected utility. Without supplied beliefs it
does not invent them.

The diagnostic can derive a first-round payoff table from the world for at
most four participants and 256 profiles. It uses the author's complete utility
specification, rather than a participant's restricted information. It does
not solve a 100-participant dynamic or incomplete-information game, and Jev
is not required to select an equilibrium. A public promise has no binding
effect unless the reviewed action rules impose one.

\subsection{Belief Updating and Repeated Worlds}\label{sec:probability}
Bayesian updating is an optional analytical calculation, distinct from action
scoring. Let $H$ denote an exhaustive set of alternative hypotheses, let $h$
be one hypothesis, and let $E$ be observed evidence. Then
\begin{equation}
P(h\mid E)=\frac{P(E\mid h)P(h)}{\sum_{h'\in H}P(E\mid h')P(h')}.
\label{eq:bayes}
\end{equation}
The prior $P(h)$ represents belief before the signal; the likelihood
$P(E\mid h)$ states how compatible that signal is with the hypothesis.
The denominator sums evidence likelihood across the alternatives so that the
updated beliefs remain normalized. Omitting it would yield unnormalized
weights rather than probabilities. A prior, likelihoods, and hypothesis set
must all be supplied; the calculation cannot derive them from a persuasive
article. Sequential updates assume conditional independence of the supplied
signals given the hypothesis. Repeated reporting of the same news can violate
that assumption and would otherwise overcount evidence.

Monte Carlo simulation means repeatedly sampling a specified model to
describe the distribution of its outcomes. Boyun can sample actions from
$q_{i,t}(a)$, while a separate event-chain model samples explicitly declared
conditional event branches. The default strategic world's external events
are scheduled, not sampled. Its current guided workflow permits only one to
four worlds; this supports path inspection, not precise frequency estimates.
Highest-score runs may repeat exactly. Increasing the number of worlds can
reduce sampling variation under a fixed model, but cannot repair incorrect
behavioral assumptions or correlated model biases.
